\section{Introduction}\label{sec:intro}

Code that runs on hardware the attacker controls can eventually be read, so the strongest
protection for an algorithm is not to ship it. Software protection has long accepted that
every client-side defence is a cost multiplier, not a barrier
\cite{collberg1997taxonomy,collberg2002watermarking,schrittwieser2016obfuscation}. Studies of
professional attackers show that they defeat layered obfuscation by combining static and
dynamic analysis and by focusing on the small region that holds the asset
\cite{ceccato2017hackers}.

Two shifts make that cost multiplier smaller. First, analysis tools are scriptable end to
end: Ghidra exposes a headless analyser and scripting API \cite{nsa2019ghidra}, frameworks
such as angr automate symbolic and dynamic analysis \cite{shoshitaishvili2016sok}, and Frida
injects JavaScript hooks into running processes to trace, patch and dump memory
\cite{frida}. Second, learned models now lift binaries toward readable source and recover
names and types: DIRE and DIRTY predict identifier names and types from decompiler output
\cite{lacomis2019dire,chen2022dirty}, and LLM4Decompile trains language models to decompile
x86 binaries end to end \cite{tan2024llm4decompile}. An automated reverse-engineering agent
(the original SA-HEACS blueprint calls such a toolkit an \enquote{REA}) can therefore chain
dumping, decompilation and semantic recovery with little human effort. The scarce resource
for the attacker is no longer analyst time; it is access to the code.

SA-HEACS removes that access for the code that matters. It splits an application into a
thin client shell, which handles input, local rendering and commodity features, and a set
of server-side services that hold the high-value IP. The idea of hiding a program slice on
a trusted server is old \cite{zhang2003hiding}, as is offloading work from a constrained
client \cite{cuervo2010maui,chun2011clonecloud,kang2017neurosurgeon}. What is new here is the
combination of three things: an explicit IP-centric partitioning rule; a learned
orchestrator that makes the economics of remote execution workable through locality
routing, micro-batching and prefetch; and a lifecycle policy that moves features between
tiers as their value decays.

This paper makes five contributions:
\begin{enumerate}
  \item An architecture that separates a zero-IP client shell from attested, segmented
        server clusters, with a request protocol (\cref{sec:arch}).
  \item A threat model that states which reverse-engineering attacks the design removes and
        which it does not (\cref{sec:threat,sec:security}).
  \item An orchestrator design with pseudocode and runnable Python for locality routing,
        micro-batching and Markov prefetch (\cref{sec:orch}).
  \item A Dynamic Feature Promotion/Demotion Pipeline that scores each feature's IP value
        and clone risk and moves it between cloud-only and client-local execution
        (\cref{sec:lifecycle}).
  \item An evaluation methodology and a reproducible simulation of cost and latency
        trade-offs, with the limits of what simulation can show (\cref{sec:eval}).
\end{enumerate}
All code is listed in \cref{app:code} and runs as shipped; every program output quoted in
the text is read from the output of that code.
